% Ad Familiares 4.1 --- headnote
% ad-familiares-04-01, around 5 April 49 BC, Laterium / Arcanum

\textit{Cicero to Servius Sulpicius Rufus, written from one of
the family properties --- the brother Quintus's estate at
Laterium or Arcanum --- around the Nones of April 49 BC,
Perseus dateline \textit{Scr. in Laterio aut Arcano Q. Ciceronis
circ. Non. Apr. a. 705 (49)}. Caesar has crossed the Rubicon
three months earlier and has now driven Pompey out of Italy and
left for Spain to face the Pompeian legions there; the consuls
and the bulk of the Senate have followed Pompey across the
Adriatic. Cicero, still in Italy and still undecided, has fallen
back from the suburbs of Rome to the countryside in Campania
and Latium. Sulpicius, consul of 51 and the most respected
jurist of the age, is in the same paralysed position --- in or
near Rome, advocating peace, ill, and uncertain whether or how
to attend the rump Senate Caesar is convening.}

\textit{The letter is the opening of a correspondence between
the two leading neutrals of the crisis. Sulpicius has asked
through their common friend G. Trebatius whether they might
meet to take counsel ``about the duty of each of us'' --- the
phrase \textit{de officio utriusque nostrum} is precisely the
Stoic technical term Cicero will work to death in De Officiis
five years later. Cicero replies in the same key: not how to
salvage anything of their old standing, but ``how to mourn the
most honourably we can.'' The reference to ``the man who was
asking me to imitate your example'' is to Caesar himself, who
had pressed Cicero to come to Rome and attend the senate
\textit{cum dignitate}; Cicero made clear he would say what
Sulpicius had said --- against the war, against the marching on
the Spains. ``A meeting of senators,'' \textit{conventus
senatorum}, is Cicero's deliberate refusal to call Caesar's
gathering a Senate.}
